\section{一个插值多项式的两种构造}
\label{sec:08-Numerical-Analysis/03-Interpolation-and-Approximation/01}

你在做实验。在 \(x=0,1,2,3\) 四个点测量了某个物理量，得到 \(y=1,3,2,5\)。你需要一个光滑的函数穿过这四个数据点，好让你估算测量点之间的值——比如 \(x=1.7\) 处大概是多少。你隐约记得多项式可以插值：\(n+1\) 个点唯一决定一个不超过 \(n\) 次的多项式。

最直接的办法是设 \(p(x)=a_0+a_1x+a_2x^2+a_3x^3\)，代入四个点得到四个方程，解 Vandermonde 线性系统。四阶系统手工可算，但三十阶的系统就会让浮点求解器叫苦。而且这还不是最要紧的——你隐约感觉插值问题一定有比硬解方程组更轻巧的办法，因为节点数据天然带有零点的结构。

\subsection{有限数据唯一决定的多项式，表达方式却不止一种}

给定 \(n+1\) 个互异节点

\[
(x_0,y_0),\ldots,(x_n,y_n),
\]

存在唯一一个次数不超过 \(n\) 的多项式 \(p_n\) 满足

\[
p_n(x_i)=y_i,\qquad i=0,\ldots,n.
\]

唯一性的证明短而有力：如果 \(p,q\) 都符合数据，那么 \(p-q\) 次数至多 \(n\)、却有 \(n+1\) 个互异零点，只能是零多项式。

这条唯一性像是法律条文——它划定边界，但不管你怎么抵达边界。同一个 \(p_n\) 可以有完全不同的表达式，这些表达式在精确数学中互等，在浮点计算中却各有软肋。Lagrange 形式和 Newton 形式描述的是同一个数学对象，却为不同的计算场景暴露不同的结构——前者让所有节点地位对称，后者方便逐点追加。

\subsection{Lagrange 形式：每个节点独占一盏灯}

Lagrange 的想法很直观：给每个节点 \(x_j\) 单独造一个基函数，让这个基函数在自己的节点处亮灯（值为 1），在别人家全部熄灭（值为 0）。然后把这些基函数的加权和当作答案——权重就是你给每个节点的 \(y_j\)。

基函数的构造一步到位：

\[
\ell_j(x)=
\prod_{m\ne j}
\frac{x-x_m}{x_j-x_m}.
\]

分子是一堆 \((x-x_m)\) 的乘积——当 \(x\) 等于除 \(x_j\) 之外的任何节点时，总有一个因子 \((x_i-x_i)=0\) 把整个乘积清零。分母是一堆 \((x_j-x_m)\) 的乘积——它是归一化常数，保证 \(\ell_j(x_j)=1\)。合起来：

\[
\ell_j(x_i)=\delta_{ij},
\]

其中 \(\delta_{ij}\) 为 Kronecker 符号：\(i=j\) 时为 1，否则为 0。插值多项式因此是

\[
p_n(x)=\sum_{j=0}^ny_j\ell_j(x).
\]

当 \(x\) 等于节点 \(x_i\) 时，除了 \(\ell_i\) 亮着，其余全部熄灭，只剩 \(y_i\)。这个构造同时给出了插值多项式的存在性证明——无需先解任何线性系统。

例如数据 \((-1,1),(0,0),(1,1)\)。三条基函数分别是 \(\ell_0(x)=x(x-1)/2\)，\(\ell_1(x)=-(x+1)(x-1)\)，\(\ell_2(x)=(x+1)x/2\)。代入求和得到 \(p(x)=x^2\)，恰好是通过这三个点的抛物线。

Lagrange 形式的优雅在于"分解再组合"——把插值问题拆解为每个节点独立贡献再叠加。但如果你要对一百万个查询点求值，每次都重新计算全部乘积，成本和舍入表现都不理想。

实际多点求值常采用 barycentric（重心）形式。预计算权重

\[
w_j=\frac{1}{\prod_{m\ne j}(x_j-x_m)},
\]

后，当 \(x\) 不等于任意节点时，插值可写为

\[
p_n(x)=
\frac{\sum_j\frac{w_jy_j}{x-x_j}}
{\sum_j\frac{w_j}{x-x_j}}.
\]

这个形式有两个好处：权重 \(w_j\) 只依赖节点位置，和函数值 \(y_j\) 无关，可以预计算并复用；求值只需 \(O(n)\) 次运算，且对 Chebyshev 型节点在数学上向前稳定。若查询点恰好等于某个 \(x_j\)，分子分母同时出现无穷，应直接返回 \(y_j\) 而非硬除。

\subsection{Newton 形式：每来一个新点，只加一层修正}

Lagrange 形式把所有节点对称地放在一起——如果你后来多加一个数据点，所有基函数都要重建。Newton 形式为增量构建而生：

\[
p_n(x)=c_0+c_1(x-x_0)
+c_2(x-x_0)(x-x_1)+\cdots
+c_n\prod_{j=0}^{n-1}(x-x_j).
\]

这个构造的妙处在于：第 \(k\) 项包含因子 \((x-x_0)(x-x_1)\cdots(x-x_{k-1})\)，在所有旧节点 \(x_0,\ldots,x_{k-1}\) 处自动为零。所以前面已经确定的项不会被后面新加的项干扰——加入第 \(k\) 项时，只需选择 \(c_k\) 来修正第 \(k\) 个新节点的值，之前节点的约束原封不动。

系数 \(c_k\) 由差商（divided difference）给出：

\[
c_k=f[x_0,\ldots,x_k],
\]

其中零阶差商就是数据本身：\(f[x_i]=y_i\)。高阶差商递归定义为

\[
f[x_i,\ldots,x_{i+k}]
=\frac{
f[x_{i+1},\ldots,x_{i+k}]
-f[x_i,\ldots,x_{i+k-1}]
}{x_{i+k}-x_i}.
\]

这个递归公式的分子是相邻两个低一阶差商的差，分母是两端节点的距离。你可以把它读作"局部斜率的斜率的斜率"——一阶差商是区间上的平均斜率，二阶差商是斜率的变化率，依此类推。

举个具体的差商表。节点 \(0,1,2\) 上的数据 \(1,2,5\)：

\[
\begin{array}{c|ccc}
x & f[\cdot] & f[\cdot,\cdot] & f[\cdot,\cdot,\cdot] \\
\hline
0 & 1 & & \\
  &   & \frac{2-1}{1-0}=1 & \\
1 & 2 &   & \frac{3-1}{2-0}=1 \\
  &   & \frac{5-2}{2-1}=3 & \\
2 & 5 & &
\end{array}
\]

系数依次为 \(c_0=1,c_1=1,c_2=1\)，于是

\[
p(x)=1+1\cdot(x-0)+1\cdot(x-0)(x-1)=1+x+x(x-1)=x^2+1.
\]

验证：\(p(0)=1,p(1)=2,p(2)=5\)。确认通过全部数据点。

Newton 形式求值可以像 Horner 法一样嵌套：从最内层括号向外，每层只做一次乘法和一次加法。加入新节点 \(x_{n+1}\) 时，差商表只需追加一个对角线元素 \(f[x_0,\ldots,x_{n+1}]\)，旧系数全部保留。这正是逐步收集数据时的理想结构。

\subsection{两种形式为什么必然殊途同归}

Lagrange 构造和 Newton 差商都产生一个次数不超过 \(n\) 且通过全部节点的多项式。由唯一性定理，它们在精确算术下是同一个 \(p_n\)。你大可以把 Newton 的差商系数展开成幂基形式，得到完全相同的结果。

但"数学上相等"和"浮点计算中表现相同"是两回事。Newton 差商的计算路径依赖节点顺序：聚集很近的节点会让差商公式中的分子（两个接近数的差）和分母（小间距）同时制造麻烦。Lagrange barycentric 形式在 Chebyshev 节点上的稳定性已经被理论证明，而幂基 Vandermonde 系统——那个一开始看起来最直白的办法——反而是条件数最糟糕的。

节点顺序不会改变差商表最终给出的数学多项式；但浮点计算中，不同的计算路径会遭受不同程度的消去。数学的不变性与实现的路径敏感性必须同时看到，这不是矛盾的——就像同一个几何事实可以用不同的坐标描述，有些坐标系在特定区域会奇异地放大微小扰动。

\subsection{插值误差：把未知函数重新带回来}

如果数据来自一个足够光滑的函数 \(f\)——而不只是一堆任意的 \((x_i,y_i)\)——我们就有了评估插值质量的基准。对每个查询点 \(x\)，插值误差为

\[
f(x)-p_n(x)
=\frac{f^{(n+1)}(\xi_x)}{(n+1)!}
\prod_{i=0}^n(x-x_i),
\]

其中 \(\xi_x\) 位于包含全部节点和查询点 \(x\) 的区间内。这个公式的结构值得多看几眼：误差被拆成两个因子的乘积。

第一个因子 \(f^{(n+1)}(\xi_x)/(n+1)!\) 完全由被插值函数的高阶导数决定。如果 \(f\) 本身就是低次多项式，这个因子直接为零——插值是精确的。如果 \(f\) 的高阶导数很大（比如在极点附近或者有快速振荡），这个因子会惩罚你。

第二个因子 \(\prod_{i=0}^n(x-x_i)\) 完全由节点的几何布局决定，和函数值无关。它是 \((n+1)\) 个距离的乘积——\(x\) 到每个节点的距离。这个乘积在节点本身上为零（误差当然为零），在节点之间则随位置变化。当 \(n\) 很大时，这个乘积在区间端点附近可以变得极大——下一节 Runge 现象的全部秘密就在这个乘积里。

增加次数 \(n\) 不一定让总误差下降，因为两个因子可能反向变化：\(f^{(n+1)}\) 可能很大，节点乘积也可能在端点爆发。增加数据点之前，先想想你要对付的函数是否真的愿意被更高次的多项式逼近。

重复节点会让普通差商的分母为零、Lagrange 基函数的分母也为零——公式在定义层面就退化了。如果同时给定了函数值和导数值，应改用 Hermite 插值，而不是把重复节点硬塞进原公式。Hermite 可看作 Newton 差商在重合节点极限下的自然推广，下一篇的数值微分也会用到类似思想。
